\documentclass[11pt,a4paper]{article}
\usepackage[margin=0.85in]{geometry}
\usepackage{hyperref}
\usepackage{enumitem}
\usepackage{graphicx}
\usepackage{xcolor}

% -----------------------------------------------------
% bvraghav-tiet-ucs503p-article.sty
% -----------------------------------------------------

% Variable: Lab Instructor
\makeatletter \newcommand{\BvrSetLabInstructor}[1]{%
  \gdef\Bvr@LabInstructor{#1}}
% default
\newcommand{\Bvr@LabInstructor}{Dr. Paramveer Sidhu}
% public accessor
\newcommand{\BvrLabInstructorValue}{\Bvr@LabInstructor}
\makeatother

% Add subtitle
\usepackage{titling}
% -- Set title bold
\pretitle{\begin{center}\bfseries\fontsize{18bp}{18bp}\selectfont}
\makeatletter
\newcommand{\BvrSubtitle}[1]{%
  \posttitle{%
    \par\end{center}
    \begin{center}\large#1\end{center}
    \vskip 0.5em}%
}
\makeatother

% Hints in the section title(s)
\newcommand{\BvrHint}[1]{%
  {\normalfont\normalsize\itshape\hfill #1}}

% Tighten list spacing to keep the proposal compact
\setlist{itemsep=0pt, topsep=2pt, parsep=0pt}
\usepackage{titlesec}
\titlespacing*{\section}{0pt}{8pt}{4pt}
\titlespacing*{\subsection}{0pt}{6pt}{3pt}

% -----------------------------------------------------

\title{SyncDrive}

\BvrSetLabInstructor{Dr. Paramveer Sidhu}

\BvrSubtitle{%
  Signal-Aware Navigation and Speed Optimization System \\
  Project Proposal for UCS503P  \\
  Submitted to: \BvrLabInstructorValue \\ \bfseries
  Thapar Institute of Engineering and Technology% 
}

\author{
  Tarun Mahindra (1024030177)
  \and 
  Parth Puri (1024030178)
  \and
  Adhiraj Kapur (1024031077)
}

\date{11 August 2026}

\begin{document}
\maketitle

\section{Higher-order goal%
  \BvrHint{\ldots secondary framing}}
SyncDrive aims to reduce avoidable time and fuel loss in
everyday urban commuting by making drivers aware of how
their speed interacts with upcoming traffic signals. The
engineering objective is to translate \emph{signal-aware
  speed guidance} into a measurable, deployable software
solution usable through a navigation application.

\section{Time-to-value %
\BvrHint{\ldots primary heuristic}}
\begin{itemize}[leftmargin=0.7cm]
    \item Deliver meaningful increments quickly by prototyping the core route-and-speed-recommendation workflow on a single simulated corridor within a short iteration window.
    \item Prioritize fast feedback over perfection using weekly iterations and early CI-backed automation.
    \item Success in iteration one is a working speed recommendation for a route with simulated signals, improved thereafter based on testing.
\end{itemize}

\section{Problem Statement}
Navigation apps such as Google Maps select routes using
distance, travel time, and traffic, but do not account for
how a driver's chosen speed interacts with signal timing.
\begin{itemize}[leftmargin=0.7cm]
    \item \textbf{Signal-blind routing:} a short, lightly congested route can still force repeated stops if driven at the speed limit.
    \item \textbf{Avoidable stop-and-go driving:} repeated braking and waiting increases journey time even in light traffic.
    \item \textbf{No dynamic speed guidance:} apps advise \emph{which road} to take, not \emph{how fast} to drive on it.
    \item \textbf{No driver-level emergency coordination:} nearby drivers get no software guidance to help clear a path for an ambulance.
\end{itemize}
\textbf{Why this matters now:} in corridors with closely spaced signals, reducing unnecessary stopping can measurably cut journey time without changing road infrastructure.

\section{Proposed Solution %
\BvrHint{\ldots Software Proposition}}
\subsection{Overview}
Build a \textbf{web/mobile-based} navigation system with:
\begin{itemize}[leftmargin=0.7cm]
    \item \textbf{Driver interface:} source/destination entry; route view with upcoming signals.
    \item \textbf{Signal timing store:} simulated/map-sourced signal location, cycle duration, expected green/red intervals.
    \item \textbf{Speed optimization engine:} computes a dynamically updated, legal recommended speed to increase the likelihood of green-phase arrivals.
    \item \textbf{Recommendation display:} e.g., ``Maintain 42 km/h for the next 800 m,'' always capped at the speed limit.
    \item \textbf{Emergency coordination module:} adjusts nearby drivers' recommendations when a registered emergency vehicle approaches.
\end{itemize}

\subsection{Core workflow}
\begin{enumerate}[leftmargin=0.7cm]
    \item Driver enters source/destination; a route is selected using distance, time, and traffic.
    \item System identifies signals along the route and retrieves simulated/known timing.
    \item Engine estimates arrival time per candidate speed and selects the legal speed minimizing expected waiting time across the route.
    \item Recommendation recalculates as position, traffic, or emergency conditions change.
\end{enumerate}

\subsection{Operational constraints}
\begin{itemize}[leftmargin=0.7cm]
    \item Recommended speed never exceeds the legal limit -- the system chooses \emph{which} legal speed, not whether to speed.
    \item No guarantee of a green signal at every intersection; guidance is expressed as increased probability and reduced expected waiting time.
    \item Real-time municipal signal-controller access is not assumed; simulated/public data is used and this limitation is stated explicitly.
\end{itemize}

\section{Solution Approach %
\BvrHint{\ldots Engineering focus}}
This is an engineering project, so we focus on a realistic, implementable design rather than novel traffic-prediction research.

\subsection{Speed optimization logic %
\BvrHint{\ldots non-research}}
For each upcoming signal:
\begin{itemize}[leftmargin=0.7cm]
    \item Determine distance from the current position to the signal.
    \item Retrieve the signal's current/predicted phase and cycle timing.
    \item Estimate arrival time for candidate legal speeds ($\text{arrival time} = \text{distance}/\text{speed}$).
    \item Estimate likely phase (green/red) and expected waiting time per candidate.
    \item Select the legal speed minimizing a weighted cost over travel time, waiting time, and stop count, considering downstream signals.
    \item Recompute whenever position, traffic, or signal data changes materially.
\end{itemize}
This is a prediction heuristic, not a guarantee of signal outcomes or a claim of novel research.

\subsection{System architecture %
\BvrHint{\ldots web-based solution}}
\begin{itemize}[leftmargin=0.7cm]
    \item \textbf{Frontend:} map, route, signal markers, recommended speed, ETA, emergency notifications.
    \item \textbf{Backend API:} route requests, signal-timing queries, optimization calculations, emergency events.
    \item \textbf{Route/traffic services:} a mapping/routing API (e.g., OpenStreetMap-based, Mapbox) for road network, route, distance, coordinates; signal timing is simulated separately since such APIs rarely expose real-time signal-phase data.
    \item \textbf{Database:} routes, simulated signal timing, requests, emergency events.
\end{itemize}

\subsection{CI/CD and fast delivery enabler}
\begin{itemize}[leftmargin=0.7cm]
    \item Build and test automatically on every change, including the optimization logic.
    \item Produce repeatable deployment artifacts to staging.
    \item Enable safe, frequent releases of incremental features (multi-signal optimization, emergency coordination).
\end{itemize}

\section{Evaluation Criterion %
\BvrHint{\ldots measurable and attributable}}
Success is measured by comparing the signal-aware strategy against a conventional constant-speed-limit strategy on identical simulated routes.

\subsection{Primary evaluation metric}
\textbf{Total journey time:} time to traverse a test route source-to-destination. Objective: reduce total journey time relative to the conventional strategy, measured experimentally. Attribution: simulated start/end timestamps under identical test conditions.

\subsection{Secondary evaluation metrics}
\begin{itemize}[leftmargin=0.7cm]
    \item Average signal waiting time and number of complete stops per route.
    \item Percentage of signals crossed during green phases.
    \item Average speed and travel distance across strategies.
    \item Emergency-vehicle travel-time improvement in the simulated coordination scenario.
\end{itemize}
No numerical improvement is assumed in advance; the prototype measures and compares these metrics experimentally.

\subsection{Pilot validation plan %
\BvrHint{\ldots high-level}}
\begin{itemize}[leftmargin=0.7cm]
    \item Define 3--5 test routes of varying length and signal density using simulated timing data.
    \item Run each route under both strategies and log all metrics.
    \item Gather informal driver feedback (5--10 participants) on clarity and usefulness.
\end{itemize}

\section{Scalability %
\BvrHint{\ldots with foresight}}
\begin{enumerate}[leftmargin=0.7cm]
    \item \textbf{Scalable (theoretically):} decoupled frontend/backend/optimization engine; indexed queries for signals/routes; stateless API design; caching of frequent route/signal data.
    \item \textbf{Operationally deployable (practically):} managed database; containerized/platform-hosted deployment; CI/CD pipeline; modular services so signal, route, and emergency data scale independently.
\end{enumerate}

\section{Engine availability heuristic}
Mature tooling is available for this project: a web/mobile frontend framework; a backend REST framework; a mapping/routing API (road network, route, distance, coordinates -- not assumed to expose real-time signal phase); geolocation services; a managed database; a simulated signal-timing dataset built for the prototype; a test framework; and a CI/CD runner. We rely on these mature components to focus on optimization-logic correctness and workflow, not on inventing infrastructure or assuming access to real municipal signal controllers.

\section{Project scope and deliverables %
\BvrHint{\ldots iteration-friendly}}
\subsection{Initial deliverable %
\BvrHint{\ldots first iteration}}
\begin{itemize}[leftmargin=0.7cm]
    \item Source/destination input and basic map interface; route generation via a routing API.
    \item Simulated traffic-signal dataset for a test corridor.
    \item Basic arrival-time calculation and speed recommendation for a single upcoming signal, capped at the speed limit.
    \item Basic backend API with unit tests for the optimization logic; CI build + tests + linting.
\end{itemize}

\subsection{Subsequent deliverables}
\begin{itemize}[leftmargin=0.7cm]
    \item Multi-signal optimization considering downstream signals.
    \item Dynamic recalculation as the vehicle moves or conditions change.
    \item Emergency-vehicle coordination module (simulated ambulance location, adjusted nearby recommendations).
    \item Dashboard comparing signal-aware vs.\ conventional strategy; performance/metric logging; scalability improvements (caching, indexing).
\end{itemize}

\section{Risks and mitigations}
\begin{itemize}[leftmargin=0.7cm]
    \item \textbf{Signal data availability:} real-time phase data is rarely public. \emph{Mitigation:} simulated/controlled data, limitation stated clearly.
    \item \textbf{Prediction uncertainty:} actual timing may differ from predicted. \emph{Mitigation:} express expected/estimated waiting time, not a guarantee; recalculate continuously.
    \item \textbf{Mapping API limitations:} usage limits or missing traffic data. \emph{Mitigation:} fall back to open/simulated signal data.
    \item \textbf{Safety:} must never encourage exceeding the speed limit. \emph{Mitigation:} engine constrains all speeds to the configured legal limit.
    \item \textbf{Emergency-coordination realism:} real infrastructure/control is inaccessible. \emph{Mitigation:} simulated emergency-vehicle environment, framed as future-integrable rather than working today.
    \item \textbf{Evaluation limitations:} a small prototype cannot reproduce city-scale traffic. \emph{Mitigation:} controlled simulations with clearly defined test scenarios.
\end{itemize}

\section{Summary}
SyncDrive proposes a deployable navigation system that recommends a legal driving speed to reduce unnecessary signal stops, alongside conventional route selection. It emphasizes time-to-value through rapid iterations, CI/CD-backed releases, and measurable evaluation (primarily total journey time, plus waiting-time and stop-count metrics), while remaining realistic about data limitations and explicitly avoiding claims of guaranteed green signals, guaranteed ambulance priority, or control of real traffic infrastructure.

\end{document}
